\documentclass[12pt]{article}
\usepackage[utf8]{inputenc}
\usepackage[T1]{fontenc}
\usepackage{hyperref}
\title{Southeast Ambrym Language Resources for K-12 Curriculum Development}
\author{Compiled by Claude Sonnet 5 (claude-sonnet-5) for Andr\'as Kornai}
\date{}
\begin{document}
\maketitle

\section{Introduction}
Southeast Ambrym (ISO 639-3 tvk; also Toak) is an Austronesian language of Ambrym, Vanuatu, with about 3,700 speakers. The dialects Endu, Penapo, Taveak and Toak are listed. The language is also called Vatlongos in some work. Documentation is mostly older academic work, and no school materials were found.

\section{Primary Educational Resources}
\begin{itemize}
\item A pedagogical study of possessive constructions in Vatlongos, written in Bislama, English and French, with a dictionary of noun possession classes: \url{https://openresearch.surrey.ac.uk/permalink/44SUR_INST/15d8lgh/alma99938165902346}
\end{itemize}

\section{Linguistic and Reference Materials}
\begin{itemize}
\item G.J. Parker, \emph{Southeast Ambrym dictionary} (Pacific Linguistics, 1970), and Parker's papers on phonology and verb inflection, listed on Nicholas Thieberger's page on the languages of Ambrym: \url{https://nthieberger.net/vanlangs/AMBRYM.HTM}
\item Glottolog entry: \url{https://glottolog.org/resource/languoid/id/sout2859}
\item Verbix verb list with glosses: \url{https://docs.verbix.com/Languages/AmbrymSoutheast}
\item Language profile in the Austronesian typology database (Toak): \url{https://lpan.eva.mpg.de/austronesian/language.php?id=1142}
\item Leibniz-ZAS Ambrym research project: \url{https://leibniz-zas.de/en/research/research-areas/project-archive/projekt-archive/ambrym}
\end{itemize}

\section{Religious and Community Resources}
None were found in the search results.

\section{Digital and Audio Resources}
\begin{itemize}
\item PARADISEC catalogue items from Ambrym: \url{https://catalog.paradisec.org.au/repository/AC2/VAMBRYM1A} and \url{https://catalog.paradisec.org.au/repository/NLA1/SPCmf01}
\item PAMBU catalogue, Ambrym Vanuatu: \url{https://asiapacific.anu.edu.au/pambu/catalogue/index.php/ambrym-vanuatu}
\end{itemize}

\section{Curriculum Development Context}
No modern reference grammar or published school materials were found. Parker's dictionary, the Vatlongos possession study and the PARADISEC recordings are the main leads. A curriculum would need new community-based orthography and reader development.
\end{document}